\documentclass[11pt]{article}
\usepackage{preamble}

\newcommand{\IconScope}{\IconSearch}

\newcommand{\valuetable}[2]{%
  \begin{center}
  \small
  \renewcommand{\arraystretch}{1.18}
  \begin{tabular}{|c|c|c|c|c|c|c|}
    \hline
    \(x\) & #1 \\ \hline
    \(y\) & #2 \\ \hline
  \end{tabular}
  \end{center}
}

\begin{document}

\packetheading{Lesson 4-1 · Period 1 · Student Packet}{Inverse Variation — Introduction}

\sectionbanner{OBJECTIVES}
{\small
\renewcommand{\arraystretch}{1.25}
\noindent\begin{tabular}{|p{1.8in}|p{4.8in}|}
\hline
\textbf{Math Objective} & Identify when two variables vary inversely by testing whether their product \(xy\) is constant, then write the inverse-variation equation \(y = k/x\) and use it to predict unknown values. \\ \hline
\textbf{Language Objective} & Students will use the frame “As \_\_\_ increases, \_\_\_ decreases; their product is \_\_\_, so the equation is \(y = \_\_\_/x\).” to explain an inverse-variation relationship. \\ \hline
\textbf{Essential Understanding} & Two variables vary inversely when their product is a constant \(k\). The equation \(y = k/x\) is the algebraic form; \(xy = k\) is the structural form. Graphically, the relationship is a hyperbola (the reciprocal function). \\ \hline
\end{tabular}
}

\sectionbanner{TOPIC GOALS / ESSENTIAL QUESTION / MATERIALS}
{\small
\renewcommand{\arraystretch}{1.25}
\noindent\begin{tabular}{|p{1.8in}|p{4.8in}|}
\hline
\textbf{Topic Goals} & Students will distinguish inverse variation from direct variation using product vs ratio tests, and write the equation \(y = k/x\) from two given values. All student-facing items are Savvas-sourced (bank IDs in teacher edition). \\ \hline
\textbf{Essential Question} & How are inverse variations related to the reciprocal function? \\ \hline
\textbf{Materials} & Pencils; student packet; calculator (optional). No Chromebooks required for Period 1. \\ \hline
\end{tabular}
}

\sectionbanner{LAUNCH — Inverse Variation (teacher models on screen)}

\begin{callout}{\IconBook\ INVERSE VARIATION RULES (keep printed on your page)}{calloutgreen}
\textbullet\ Two variables vary INVERSELY when their product is a constant \(k\):\\
\hspace*{1.6em}\(xy = k\) \hspace{1.5em} or equivalently \hspace{1.5em} \(y = k/x\)\\
\textbullet\ The constant \(k\) is the CONSTANT OF VARIATION.\\
\textbullet\ As one variable INCREASES, the other DECREASES proportionally.\\
\textbullet\ Example: if \(x = 10\) and \(y = 3\) vary inversely, then \(k = xy = 30\).\\
\hspace*{1.6em}Equation: \(y = 30/x\). When \(x = -6\): \(y = 30/(-6) = -5\).
\end{callout}

\begin{callout}{\IconWarn\ COMMON ERROR}{calloutred}
Do NOT write \(y/x = k\) — that is DIRECT variation.\\
Inverse variation multiplies, direct variation divides.\\
Test: compute \(xy\) for several points. If constant \(\rightarrow\) inverse.
\end{callout}

\sectionbanner{EXPLORE — Think-Pair-Share (33 min)}
\textit{Work with a partner. Use the Inverse Variation Rules above for every problem. Move through items in order — each builds on the last.}

\textbf{Bank-item labels (in order):}
\begin{enumerate}[leftmargin=1.6em,itemsep=2pt,topsep=2pt]
  \item Try It 1 — Identify inverse variation
  \item Practice \#9 — Generalize
  \item Practice \#14 — Another table check
  \item Try It 2 — Write the equation
  \item Practice \#13 — Generalize: write \(k\) in terms of \(x\) and \(y\)
  \item Practice \#16 — Apply the equation
  \item Practice \#10 — Construct Arguments (extension)
\end{enumerate}

\bankitem{Try It 1 — Identify inverse variation}{%
Try It 1. Determine if each table of values represents an inverse variation. (a) \(x\): 1,2,3,5,6,15; \(y\): 25.5, 12.75, 8.50, 5.10, 4.25, 1.70. (b) \(x\): 6.6, 5.5, 4.4, 3.3, 2.2, 1.1; \(y\): 3, 5, 7, 9, 11, 13.

\medskip
\textbf{(a)}
\valuetable{1 & 2 & 3 & 5 & 6 & 15}{25.5 & 12.75 & 8.50 & 5.10 & 4.25 & 1.70}

\textbf{(b)}
\valuetable{6.6 & 5.5 & 4.4 & 3.3 & 2.2 & 1.1}{3 & 5 & 7 & 9 & 11 & 13}
}{}

\bankitem{Practice \#9 — Generalize}{%
9. Generalize: Just from looking at the table of values, how can you determine that the data do NOT represent an inverse variation? Table - \(x\): -2, 2, 4, 6, 8, 10; \(y\): -6, 6, 12, 18, 24, 30.

\valuetable{-2 & 2 & 4 & 6 & 8 & 10}{-6 & 6 & 12 & 18 & 24 & 30}
}{}

\bankitem{Practice \#14 — Another table check}{%
14. Do the tables of values represent inverse variations? Explain. See Example 1. Table - \(x\): -1/4, -1/2, 1/3, 2, 5, 11; \(y\): -9/2, -9, 6, 36, 90, 198.

\valuetable{\(-\frac{1}{4}\) & \(-\frac{1}{2}\) & \(\frac{1}{3}\) & 2 & 5 & 11}{\(-\frac{9}{2}\) & -9 & 6 & 36 & 90 & 198}
}{}

\bankitem{Try It 2 — Write the equation}{%
Try It 2. In an inverse variation, \(x = 6\) when \(y = \frac{1}{2}\). (a) What is the equation that represents the inverse variation? (b) What is the value of \(y\) when \(x = 15\)?
}{}

\bankitem{Practice \#13 — Generalize: write \(k\) in terms of \(x\) and \(y\)}{%
13. Generalize: For an inverse variation, write an equation that gives the value of \(k\) in terms of \(x\) and \(y\).
}{}

\bankitem{Practice \#16 — Apply the equation}{%
16. If \(x\) and \(y\) vary inversely and \(x = 3\) when \(y = \frac{2}{3}\), what is the value of \(y\) when \(x = -1\)?
}{}

\bankitem{Practice \#10 — Construct Arguments (extension)}{%
10. Construct Arguments: Explain why zero cannot be in the domain of an inverse variation.
}{}

\sectionbanner{SHARE / SUMMARY (5 min)}

\begin{sentenceframebox}
\textbf{Sentence frame}

Pair share: “As \_\_\_ increases, \_\_\_ decreases. Their product is \_\_\_, so the equation is \(y = \_\_\_/x\).”

Whole class: one pair reads their sentence. Class signals thumbs up/sideways.
\end{sentenceframebox}

\sectionbanner{EXIT}

\begin{summaryexitbox}{EXIT TICKET — Summary Recap}
Today I learned that \_\_\_\_\_\_\_\_\_\_\_\_\_\_\_\_\_\_\_\_\_\_\_\_\_\_\_\_\_\_\_\_\_\_\_\_\_\_\_\_\_\_\_\_\_\_\_\_.

The difference between inverse and direct variation is \_\_\_\_\_\_\_\_\_\_\_\_\_\_\_\_\_\_\_\_\_\_\_\_\_\_\_\_\_\_\_\_\_\_\_\_\_\_\_\_\_\_\_\_\_\_\_\_.

One thing I'm still unsure about is \_\_\_\_\_\_\_\_\_\_\_\_\_\_\_\_\_\_\_\_\_\_\_\_\_\_\_\_\_\_\_\_\_\_\_\_\_\_\_\_\_\_\_\_\_\_\_\_.
\end{summaryexitbox}

\clearpage

\sectionbanner{OPTIONAL REINFORCEMENT (not collected)}
\textit{If you finish Explore early or want extra practice before Period 2.}

\begin{callout}{\IconScope\ OPTIONAL · Model \& Discuss Parts B + C}{calloutblue}
Return to the two rectangles from the Do Now (both have area 144). With your partner:

\medskip
B) Use Structure. Considering rectangles with area 144, what happens to the WIDTH as the LENGTH increases?

\medskip
C) Examine at least five more pairs of rectangles with area 144. How would you describe the relationship between their lengths and widths — in your own words?
\end{callout}

\bankitem{Optional \#1 (Savvas Practice)}{%
15. Do the tables of values represent inverse variations? Explain. See Example 1. Table - \(x\): 1, 2, 3, 4, 5, 6; \(y\): 60, 30, 20, 15, 12, 10.

\valuetable{1 & 2 & 3 & 4 & 5 & 6}{60 & 30 & 20 & 15 & 12 & 10}
}{}

\end{document}
